\scanpage{208}
% archive sheet number: 207

447 \quad ($^{\ast\handwritten{\ast}}$) Les deux autres livres sont les thèses de Jean Giraud et de \deleted{Monique} Monique Hakim (sur le formalisme des champs et de la \unsure{1|l}-cohomologie non commutative, et sur les schémas relatifs sur des topos annelés généraux).

448 \quad ($^{\ast}$) Cette citation (traduite) est extraite de la courte lettre (adressée au Dr. Peters) du 18 \unsure{février|févriier} l'an dernier, où je l'informais de ma décision de me retirer de l'``editorial board'' des Grundlehren. Le Dr. Peters avait en fait quitté déjà le Springer Verlag (il travaille à présent dans le Birkhäuser Verlag), et la correspondance s'est continuée avec le \deleted{\unsure{Dr.}} Dr. J. Heinze, en charge des Grundlehren dans la maison Springer\handwritten{.} \deleted{Verlag.} J'avais demandé\deleted{\unsure{e}} qu'une copie de ma lettre soit envoyée à chacun des co-éditeurs des Grundlehren (au nombre de dix-huit), et \unsure{avais|avaàs} réitéré cette prière au Dr Heinze\deleted{\unsure{,}} à deux reprises (en avril 84 et janvier 85) sans qu'il juge utile de me \deleted{dire} préciser si oui ou non elle avait été respectée (il \inserted{\handwritten{est}} appar\deleted{\unsure{aît|ait}}\handwritten{\unsure{ent|ut}} que \uline{non}). J'ai pris la peine d'envoyer moi-même une copie de \unsure{ma|la} lettre \deleted{de démission, avec quelques mots d'explication,} à chacun des dix-huit éditeurs, avec quelques mots d'explication \deleted{adressé} au sujet de cet envoi. \deleted{Parmi ces éditeurs, j'en conn} \inserted{Je} connais bien personnellement sept parmi \deleted{\unsure{ces}}\handwritten{eux}\handwritten{,} \deleted{éditeurs,} et \deleted{cinq entre eux} \handwritten{j'en} comptais \inserted{\handwritten{cinq}} parmi mes amis. Un seul \deleted{\unsure{(}} (Artin) a pris la peine de me répondre, et aucun apparemment \deleted{\unsure{n'a}}\handwritten{n'a} rien trouvé d'anormal \deleted{visà} (ne serait-ce que vis-à-vis de \mbox{l\,u\,i\,-\,m\,ê\deleted{\unsure{e}}m\,e} ) que la maison Springer n'ait pas \inserted{\handwritten{elle-même}} pris la peine de lui faire parvenir \deleted{elle-même} \deleted{au mois} (et dès le mois de février 1984) la lettre en question.

($^{\ast\ast\handwritten{\ast}}$) Ce sont les volumes n\textsuperscript{\uline{o}}s 288,340, 569, 589, 900.

($^{\ast\ast\ast\handwritten{\ast}}$) Comme je le précise dans l'avant-dernière note de b. de p. , le Dr. Peters a depuis quitté \inserted{\handwritten{le}} Springer Verlag pour Birkhäuser Verlag.

450 \quad ($^{\ast}$) \deleted{Voir mes observations \inserted{\unsure{(p.899)}} dans la sous-note ``Les double-sens - ou l'art de l'arnaque'' (n\textsuperscript{\uline{o}} 169$_7$), concernant} \inserted{\handwritten{Dans}} le fameux ``Fil d'Ariane'' (\deleted{à travers mon}\inserted{\handwritten{à travers SGA 4\unsure{, etc| et 5}}}) dans le volume nommé `` SGA $4\frac12$ '', \deleted{\unsure{on}} rien ne pourrait laisser supposer au lecteur que j'aie eu l'honneur de faire des exposés dans SGA 4 et SGA 5 (par contre, j'ai \inserted{\handwritten{bien}} eu celui de ``collaborer'' à ``SGA $4\frac12$ '' ...). \inserted{\handwritten{Voir à ce sujet mes observations dans la note ``Les double-sens - ou l'art de l'arnaque'' (n\textsuperscript{o} $169_{\unsure{7|2}}$), p. 899.}}

\unsure{451|452} \quad ($^{\ast}$) J'ai songé d'ailleurs aussi qu'il se pourrait bien que la situation se renverse, et que ce soit l'estimable entre\deleted{\unsure{m|p}}prise qui m'intente un procès, pour atteinte à sa réputation. \deleted{J'imagine que} \deleted{\unsure{c}}\handwritten{C}es gens ``au service de la \deleted{\unsure{s}}\handwritten{S}cience'', ils doivent être pointilleux sur ce chapitre \deleted{\unsure{en}} (du moment que c'est de \uline{leur} \deleted{propre} réputation qu'il s'agit ...).

452 \quad ($^{\ast}$) Cette idée-force de bâtir des ``maisons'', et qui soient bonnes ``pour tous'', a joué un rôle considérable dans mon oeuvre mathématique, et ceci depuis le début des années cinquante déjà. Cela a été l'expression concrète \deleted{de la ``pulsion de} \inserted{\handwritten{dans mon travail de}} ce que j'ai appelé la ``pulsion de service'' \deleted{en moi}, qui a fait partie (sans même que je la décèle avant la réflexion ``La clef du yin et du yang'') des \deleted{motivations profon} forces profondes donnant leur force vive à mon travail mathématique. L'archétype de la ``maison'' apparaît pour la première fois dans ma réflexion, \deleted{et avec} sans que je \inserted{ne} l'aie \deleted{entrev} pressent\inserted{\unsure{e}}\unsure{i}, et avec une grande force, dans la note \inserted{du 26 novembre} ``Yin le serviteur, et les nouveaux maîtres'' (n\textsuperscript{\uline{o}} 135).
